% Chapter 3 -- Working with Files and Directories.
% Sources (Law 1): os/milton/command.c (builtin_ren AH=56h same directory;
% builtin_cd; builtin_del wildcards); os/milton/int21.c (8.3 search template
% '*' '?', do_chdir: relative, absolute, ".."; "\" alone is a no-op in the
% shipped build -- bead needed, WL-0119); os/milton/dosutil.c (ATTRIB
% wildcard display); os/milton/xcopy_tree.c (TREE).
% GATES: test-command test-ut6d test-fat-subdir test-attrib test-xcopy-tree test-pubs-screens
\Gates{test-command, test-ut6d, test-fat-subdir, test-attrib, test-xcopy-tree, test-pubs-screens}
\chapter{Working with Files and Directories}

\section{Naming Files}
\index{file names}
Chapter~1 described the two parts of a file name: a name of up to eight
characters and an extension of up to three.  Choose names that tell you
what a file holds, and use the extension to group files of one kind:
\file{.TXT} for text, \file{.BAK} for backup copies, \file{.BAT} for batch
files.  Make your names from letters and digits; do not use spaces, or the
characters \texttt{\textbackslash{} / : * ? " < > | + = ; , [ ]}, which have
other meanings on the command line.

\section{Renaming a File}
\index{renaming files}\index{REN@\texttt{REN}}
The \cmd{REN} command gives a file a new name.  \cmd{RENAME} is another
name for the same command.  Name the file first and the new name second.
\begin{Steps}
\item Type \type{CLS} and press \key{Enter}.
\item Type \type{REN MEMO.TXT MEMO.OLD} and press \key{Enter}.
\item Type \type{DIR} and press \key{Enter}.
\end{Steps}

\Screen[rows=16, label=fig:ren]{screens/ren.txt}{Renaming MEMO.TXT}

The file keeps its size, date and directory; only the name changes.  If a
file with the new name already exists, or the file you named does not,
\cmd{REN} displays \texttt{Bad command or file name} and nothing is
changed.  To put the old name back, type \type{REN MEMO.OLD MEMO.TXT}.

\section{Naming Groups of Files}
\index{wildcards}
Some commands can act on a group of files at once.  To name the group, use
the \term{wildcard} characters:

\SideHead{* (asterisk)}Stands for any characters from its position to the
end of the name, or of the extension.  \file{*.TXT} means every file with
the extension \file{TXT}; \file{MEMO.*} means every file named \file{MEMO},
whatever its extension; \file{*.*} means every file.

\SideHead{? (question mark)}Stands for any one character in its position.
\file{MEMO.?AK} matches \file{MEMO.BAK} but not \file{MEMO.TXT}.

The commands that accept wildcards are \cmd{DEL}, \cmd{ATTRIB},
\cmd{CHKDSK} and \cmd{XCOPY}; the \doctitle{InitechDOS Reference} states it
for each command.  To see which files a wildcard name matches, give it to
\cmd{ATTRIB}, which lists every match:
\begin{Steps}
\item Type \type{CLS} and press \key{Enter}.
\item Type \type{ATTRIB *.TXT} and press \key{Enter}.
\end{Steps}

\Screen[rows=5, label=fig:wild]{screens/wild.txt}{The files that match *.TXT}

\begin{Important}
Check a wildcard name with \cmd{ATTRIB} before you give it to \cmd{DEL}:
\cmd{DEL} deletes every file the name matches.
\end{Important}

\section{The Current Directory}
\index{current directory}
The directory you are working in is the \term{current directory}.  Commands
look for files in the current directory unless you tell them otherwise, and
the command prompt always shows it.  The practice diskette has two
subdirectories, \file{APPS} and \file{SAMPLES}.  In this exercise you move
into \file{SAMPLES} and back:
\begin{Steps}
\item Type \type{CLS} and press \key{Enter}.
\item Type \type{CD SAMPLES} and press \key{Enter}.  The prompt becomes
  \file{A:\bs SAMPLES>}.
\item Type \type{CD} and press \key{Enter}.  \cmd{CD} alone displays the
  current directory.
\item Type \type{CD ..} and press \key{Enter} to go back to the parent
  directory, the root.
\item Type \type{CD}, then \type{CD SAMPLES}, then \type{DIR}, pressing
  \key{Enter} after each.
\end{Steps}

\Screen[rows=25, label=fig:cdsub]{screens/cdsub.txt}{Moving between directories}

\section{Paths}
\index{paths}
To name a file in a directory other than the current one, write its
\term{path}: the names of the directories that lead to it, each followed by
a backslash (\type{\bs}).
\begin{itemize}
\item A path that begins with a backslash starts at the root directory:
  \file{\bs SAMPLES\bs GREET.BAT} is the file \file{GREET.BAT} in the
  directory \file{SAMPLES}, which is in the root directory.
\item Any other path starts at the current directory: from the root,
  \file{SAMPLES\bs GREET.BAT} names the same file.
\item \file{..} names the parent directory: from \file{SAMPLES},
  \file{..\bs MEMO.TXT} is the file \file{MEMO.TXT} in the root.
\end{itemize}
You may put the drive in front of a path: \file{A:\bs SAMPLES\bs GREET.BAT}.

\draft{bead needed: CD \bs{} alone leaves the current directory unchanged in
the shipped build (int21.c, do\_chdir); DOS 3.3 returns to the root.  Use
CD .. until then}
To go straight to the root directory from anywhere on the disk, type
\type{CD \bs}.

\section{The Directory Tree}
\index{directory tree}\index{TREE@\texttt{TREE}}
Because directories can hold other directories, the directories of a disk
form a \term{tree}, with the root directory at its base.  The \cmd{TREE}
utility program displays it:
\begin{Steps}
\item Type \type{CLS} and press \key{Enter}.
\item Type \type{TREE} and press \key{Enter}.
\end{Steps}

\Screen[rows=16, label=fig:tree]{screens/tree.txt}{The directories of the practice diskette}

For each directory \cmd{TREE} shows its path and the subdirectories it
holds.  Type \type{TREE /F} to see the files in each directory as well.

\section{Organizing Your Files}
A few habits keep a disk easy to use:
\begin{itemize}
\item Keep programs in one directory and your own files in others, one
  directory to a subject or a project.
\item Keep the root directory short: the practice diskette's root holds
  only a few files and the two directories.
\item Remove files you no longer need, and directories that are empty.
\item Make copies of important files on another diskette regularly
  (Chapter~4).
\end{itemize}
